\section{Reproductibilité des figures et des tableaux}
\label{sec:S7}
\sloppy


Cette partie indique, pour chaque figure tirée des données, comment la reproduire à partir des fichiers déposés avec l'article et de l'application Isomorph-Reborn.

\subsection*{Réseau ISOMORPH (figure~\ref{fig:d2})}

Dans l'application Isomorph-Reborn, fournie avec l'article, ouvrir l'onglet « Générateur de scénarios » ; dans la section « Réseau », en tête de l'onglet, choisir « Réseau ISOMORPH » pour afficher le schéma par échelons de la figure~\ref{fig:d2}. La description du réseau, avec ses capacités et ses délais de trajet, est contenue dans le fichier \texttt{reseau\_isomorph-originel.json}.

\subsection*{Crise S0366 sous les quatre politiques (figure~2 du manuscrit)}

Dans l'onglet « Générateur de scénarios » d'Isomorph-Reborn, charger le fichier d'hypothèses \texttt{hypotheses\_lot} \texttt{\_A\_isomorph\_tension75.json} (réseau ISOMORPH, représentation par flux, niveau de tension 75, 2 000 scénarios, graine 42) et lancer la génération. Les trajectoires journalières sont exportées dans le fichier \texttt{ip\_timeseries.csv} (colonnes \texttt{scenario\_id}, \texttt{branch\_id}, \texttt{day\_rel} et \texttt{ip}) ; la figure~2 du manuscrit en reprend les quatre versions du scénario S0366. Les jours de détection proviennent de la base d'apprentissage exportée pour le même lot.

\subsection*{Ajustement de la crise S0366 (figure~\ref{fig:d4})}

Les paramètres ajustés de chaque version de crise sont exportés par Isomorph-Reborn (onglet « Ajustement par lot ») dans le fichier \texttt{resilience\_ip.csv} (colonnes \texttt{k}, \texttt{q}, \texttt{h}, \texttt{g}, \texttt{d}, \texttt{alpha}, \texttt{beta}, \texttt{r2}, puis les six indicateurs). La courbe tracée est l'équation (2) du manuscrit évaluée avec les paramètres de la ligne S0366, version \texttt{none}, superposée à la trajectoire du fichier \texttt{ip\_timeseries.csv}. Dans l'application, cette crise est la courbe n° 1461 sur 8 000 : le compteur numérote les versions de crise, à raison de quatre par scénario, et non les scénarios.

\subsection*{Apprentissage du réseau bayésien (figure~4 du manuscrit)}

Dans BayesArtist, importer la base \texttt{bayesartist\_apprentissage\_lotA.csv} (6 400 versions de crise, 20 variables) et les contraintes \texttt{bayesartist\_contraintes\_isomorph.json} (155 arcs interdits, quatre parents au plus), puis lancer l'apprentissage par recherche gloutonne (\textit{Greedy Hill Climbing}) au score BIC avec lissage de Laplace. Le réseau appris s'exporte au format \texttt{.bif} ; la figure est tracée à partir de ce fichier. BayesArtist n'étant pas diffusé, le même apprentissage peut être conduit avec tout logiciel de réseaux bayésiens qui apprend à partir d'un fichier CSV et accepte des arcs interdits (voir la déclaration de disponibilité des données du manuscrit).

\subsection*{Pronostic de l'étude de cas (figure~\ref{fig:d7})}

Dans BayesArtist, en mode inférence sur le réseau de la figure~4 du manuscrit, poser les évidences nature = \texttt{fermeture\_entrepot}, cible = \texttt{passage\_oblige} et durée = \texttt{fort}, puis fixer successivement le profil à \texttt{none}, \texttt{tardif}, \texttt{prudent} et \texttt{reactif}, et lire à chaque fois la distribution du nœud IPC.

\subsection*{Carte de préconisation (figure~5 du manuscrit)}

Dans BayesArtist, sur le réseau de la figure~4 du manuscrit, poser les évidences nature, cible et durée correspondant à un contexte de la carte, ainsi que IPC = \texttt{faible}, puis lire la distribution du nœud profil. La carte est calculée par inférence exacte sur le réseau exporté au format \texttt{.bif} ; BayesArtist, dont l'inférence est approchée, donne les mêmes valeurs à 0,01 près (par exemple 0,06, 0,08, 0,10 et 0,76 pour la fermeture longue d'un point de passage obligé).

\subsection*{Lecture économique (tableau~9 et figure~6 du manuscrit)}

Les ventes perdues sont calculées à partir du fichier \texttt{ip\_timeseries.csv} du lot A, comme la somme de $1 - \mathrm{IP}(t)$ sur les 60 jours suivant le choc multipliée par la demande de référence ; le nombre de leviers engagés et la demande de référence proviennent de la base d'apprentissage exportée pour le même lot (colonnes \texttt{nb\_decisions} et \texttt{base\_demand}).
